\section{Introduction}
\label{sec:introduction}

This is the fourth presentation of Vibe Theory and the most complete. The first
(Pollard, 2025) set out the framework as philosophy. The second (Pollard, 2026a) gave
it a discrete mathematical structure. The third (Pollard, 2026b) tested the conceptual
core in simulation. This paper does two things at once. It states the model in its current form, and it walks through the
full set of findings, thirty-six in all, saying for each what was found and where we
landed.

The thesis is monism. The framework treats reality as a single kind of thing, the vibe,
a unit of experience, whose only act is to experience other vibes and be experienced by
them. In the model there is no background space or time. Space is then the shape of the
web of vibes seen from far away, time is the depth of its before-and-after, and
everything else, matter, force, the quantum, is modeled as a pattern in the web. A monist theory of this strictness owes a debt
most ontologies never pay: it must show that the variety of the physical world actually
comes out of one substrate, not merely assert that it could. This paper pays that debt
as far as a finite simulation can.

The method is to turn each load-bearing question into a runnable measurement. We built
a testbed, a finite and seeded program with a known-answer test suite, and posed the
questions the framework must answer to be coherent, each as an experiment that returns
a number and that could fail. Several early versions did fail, in plausible-looking
ways, and were caught by the known-answer tests and corrected. What survives is
reported here, sorted honestly into what is demonstrated, what is a first rung on a deep
problem, and what is open or beyond reach.

The paper is organised as follows. Section~\ref{sec:model} states the current model, the substance, the relation, the rule, and the growth, with the small constructor
that writes it at a glance. Section~\ref{sec:testbed} describes the testbed and the
known-answer discipline that makes its numbers trustworthy. The result sections then
walk through the findings by theme: foundations (geometry, time, the law, the dominant
spacetime phase), Lorentz invariance, matter and the fields, forces and gravity, the
quantum sector, cosmology and the dark sector, holography, the distinctive predictions
and their contact with data, and the capstone where the whole model is run as one
system. Section~\ref{sec:scoreboard} is the full scoreboard of all thirty-six findings
with their honest status.

Throughout, the framework sits in the family of discrete-physics programs, the causal
sets of Bombelli, Lee, Meyer, and Sorkin (1987), the lattice gauge theory of Wilson
(1974), the overlap fermions of Neuberger (1998), and the cellular-automaton
interpretation of quantum mechanics of 't Hooft (2016). What Vibe Theory adds is the
insistence that all of these live on one experiential substrate, and the demonstration,
here, that they can. This work is a philosophical and mathematical model with a
simulable build target rather than a finished scientific theory. References to physics
are sources of structure and tests of coherence, not claims to have derived nature.
